% !TeX root = main.tex

% Аргументы: имя файла в images, подпись, метка для ссылки.
% После добавления файла рамка автоматически заменяется рисунком.
\newcommand{\imageplaceholder}[3]{%
  \begin{figure}[!htbp]
    \centering
    \IfFileExists{images/#1}{%
      \includegraphics[width=0.95\textwidth,keepaspectratio]{#1}%
    }{%
      \fbox{%
        \parbox[c][3cm][c]{0.9\textwidth}{%
          \centering
          Место для рисунка\\[0.4em]
          \texttt{images/\detokenize{#1}}
        }%
      }%
    }
    \caption{#2}
    \label{#3}
  \end{figure}
}

% Заменять вызовы команды готовым текстом по мере выполнения работы.
\newcommand{\reporttodo}[1]{%
  \par\noindent\textit{[Заполнить: #1.]}\par
}

\tableofcontents
\newpage

\section{Цель работы}

Освоить принцип адаптивного слежения за эталонным сигналом
для неустойчивого многомерного линейного объекта. Исследовать влияние
коэффициента адаптации на ошибку слежения и управляющее воздействие.

\section{Постановка задачи и исходные данные}

Рассматривается объект управления
\begin{equation}
  \dot{\bm x}=A\bm x+\bm b u,
  \qquad y=C\bm x,
  \qquad \bm x\in\mathbb R^2.
  \label{eq:plant}
\end{equation}
Вектор состояния $\bm x(t)$ и выход $y(t)$ доступны измерению,
матрицы объекта известны.

Заданы
\begin{equation}
  A=\begin{pmatrix}3&1\\1&1\end{pmatrix},
  \qquad
  \bm b=\begin{pmatrix}1\\2\end{pmatrix},
  \qquad t_{\text{п}}=0{,}16\,\text{с},
  \qquad \sigma=0\%.
  \label{eq:variant}
\end{equation}
\reporttodo{указать матрицу выхода $C$ и соответствующее выражение для $y$}

Используются задающее воздействие и коэффициенты полинома фильтра:
\begin{equation}
  g(t)=\sin t,
  \qquad k_{g1}=0{,}8,
  \qquad k_{g0}=0{,}16.
  \label{eq:reference-data}
\end{equation}
Сигнал $g(t)$ доступен измерению, однако амплитуда, частота
и фаза его гармоник априори неизвестны регулятору.

Требуется построить управление, обеспечивающее ограниченность
всех сигналов замкнутой системы и выполнение условия
\begin{equation}
  \varepsilon(t)=g(t)-y(t),
  \qquad \lim_{t\to\infty}\varepsilon(t)=0.
  \label{eq:goal}
\end{equation}

\section{Синтез стабилизирующей обратной связи}

\subsection{Проверка управляемости объекта}

Матрица управляемости имеет вид
\begin{equation}
  \mathcal C=\begin{pmatrix}\bm b&A\bm b\end{pmatrix}.
  \label{eq:controllability}
\end{equation}
\reporttodo{подставить матрицы объекта, вычислить определитель
и ранг $\mathcal C$, сделать вывод об управляемости}

\subsection{Выбор желаемой динамики}

Для апериодического процесса рассматривается стандартный полином
\begin{equation}
  D_{\text{ж}}(s)=(s+\omega_0)^2,
  \qquad
  A_{\text{ж}}=
  \begin{pmatrix}0&1\\-\omega_0^2&-2\omega_0\end{pmatrix},
  \qquad \omega_0=\frac{t_{\text{п}}^*}{t_{\text{п}}}.
  \label{eq:desired-dynamics}
\end{equation}
\reporttodo{указать полосу допуска для определения времени переходного
процесса, найти $t_{\text{п}}^*$, рассчитать $\omega_0$
и численные элементы $A_{\text{ж}}$}

\subsection{Расчёт матрицы обратной связи}

Матрица $K$ определяется методом модального управления:
\begin{equation}
  AM-MA_{\text{ж}}=\bm b H,
  \qquad K=HM^{-1},
  \qquad A_M=A-\bm b K.
  \label{eq:modal-control}
\end{equation}
Матрица $H$ выбирается так, чтобы пара $(A_{\text{ж}},H)$
была полностью наблюдаемой.

\reporttodo{выбрать $H$, привести решение уравнения Сильвестра $M$,
проверить его невырожденность, записать $K$ и $A_M$;
проверить собственные значения $A_M$}

\noindent\textbf{Вывод по пункту:}
\reporttodo{сформулировать результат проверки управляемости
и синтеза стабилизирующей обратной связи}

\section{Формирование регрессора задающего сигнала}

\subsection{Модель генератора и параметризация сигнала}

Для сигнала~\eqref{eq:reference-data} модель генератора имеет вид
\begin{equation}
  \ddot g+l_1\dot g+l_0g=0,
  \qquad l_1=0,
  \qquad l_0=1.
  \label{eq:reference-generator}
\end{equation}
Параметризованное представление сигнала:
\begin{equation}
  g=\bm\theta_g^{\mathsf T}\bm\xi_g+\delta_g(t),
  \qquad
  \bm\theta_g=\begin{pmatrix}k_{g0}-l_0\\k_{g1}-l_1\end{pmatrix},
  \qquad \lim_{t\to\infty}\delta_g(t)=0,
  \label{eq:reference-parameterization}
\end{equation}
где $\delta_g(t)$ учитывает рассогласование начальных условий
генератора и фильтра.

\reporttodo{записать численный вектор $\bm\theta_g$
и пояснить роль начальных условий в параметризации}

\subsection{Построение фильтра}

Регрессор $\bm\xi_g$ формируется фильтром
\begin{equation}
  \dot{\bm\xi}_g=A_{0g}\bm\xi_g+\bm b_{0g}g,
  \qquad
  \bm\xi_g=\begin{pmatrix}\xi_{g1}\\\xi_{g2}\end{pmatrix},
  \label{eq:reference-filter}
\end{equation}
где
\begin{equation}
  K_g(s)=s^2+k_{g1}s+k_{g0},
  \qquad
  A_{0g}=\begin{pmatrix}0&1\\-k_{g0}&-k_{g1}\end{pmatrix},
  \qquad \bm b_{0g}=\begin{pmatrix}0\\1\end{pmatrix}.
  \label{eq:filter-matrices}
\end{equation}
\reporttodo{подставить коэффициенты фильтра, проверить гурвицевость
$A_{0g}$ и привести скалярные уравнения реализации}

\noindent\textbf{Вывод по пункту:}
\reporttodo{указать, какие сигналы использует фильтр
и почему его выходы ограничены}

\section{Адаптивный следящий регулятор}

\subsection{Закон управления и модель ошибки}

Управление складывается из стабилизирующей и адаптивной частей:
\begin{equation}
  u=-K\bm x+u_a,
  \qquad u_a=\hat{\bm\psi}_g^{\mathsf T}\bm\xi_g,
  \label{eq:control}
\end{equation}
где $\hat{\bm\psi}_g$ является вектором настраиваемых параметров.
Передаточная функция стабилизированной части системы:
\begin{equation}
  W(s)=C(sI_2-A_M)^{-1}\bm b.
  \label{eq:stable-transfer}
\end{equation}
\reporttodo{вычислить $W(s)$ для выбранного выхода,
записать модель ошибки слежения и условия её применимости}

\subsection{Расширенная ошибка и алгоритм адаптации}

Введём фильтрованные сигналы
\begin{equation}
  \bm\omega=W(s)[\bm\xi_g],
  \qquad \nu=W(s)[u_a].
  \label{eq:filtered-signals}
\end{equation}
Оператор $W(s)$ применяется к каждой компоненте $\bm\xi_g$.
Расширенная ошибка и алгоритм настройки имеют вид
\begin{equation}
  \hat\varepsilon
    =\varepsilon-\hat{\bm\psi}_g^{\mathsf T}\bm\omega+\nu,
  \qquad
  \dot{\hat{\bm\psi}}_g=\gamma\bm\omega\hat\varepsilon,
  \qquad \gamma>0.
  \label{eq:adaptation}
\end{equation}
\reporttodo{записать реализацию $W(s)$; пояснить смысл
расширенной ошибки и свойства алгоритма
с учётом начальных переходных процессов}

\clearpage
\subsection{Реализация и условия моделирования}

\reporttodo{указать начальные состояния объекта и фильтров,
начальную оценку $\hat{\bm\psi}_g(0)$, интервал моделирования,
решатель, шаг или допуски интегрирования}

Для сравнения выбираются два различных коэффициента адаптации:
\begin{equation}
  \gamma_1=\ldots,\qquad \gamma_2=\ldots,
  \qquad \gamma_1\ne\gamma_2.
  \label{eq:adaptation-gains}
\end{equation}
В обоих опытах сохраняются одинаковые начальные условия,
задающее воздействие, матрицы обратной связи и параметры фильтров.
Схема моделирования представлена на рисунке~\ref{fig:scheme}.

\imageplaceholder{lab12_scheme.png}
  {Схема адаптивной системы слежения в Simulink}
  {fig:scheme}

\clearpage
\section{Исследование системы при первом коэффициенте адаптации}

В первом опыте используется коэффициент $\gamma=\gamma_1$
из~\eqref{eq:adaptation-gains}.
На рисунках~\ref{fig:gamma1-error} и~\ref{fig:gamma1-control}
представлены ошибка слежения и управляющее воздействие.

\imageplaceholder{lab12_gamma1_error.png}
  {Ошибка слежения $\varepsilon(t)$ при $\gamma=\gamma_1$}
  {fig:gamma1-error}

\imageplaceholder{lab12_gamma1_u.png}
  {Управляющее воздействие $u(t)$ при $\gamma=\gamma_1$}
  {fig:gamma1-control}

\subsection{Анализ результатов}

\reporttodo{описать изменение ошибки слежения, оценить время
её затухания и остаточные колебания, отметить максимальное управление
и характер его установившегося режима}

\noindent\textbf{Вывод по пункту:}
\reporttodo{оценить достижение цели слежения
при первом коэффициенте адаптации по результатам моделирования}

\clearpage
\section{Исследование системы при втором коэффициенте адаптации}

Во втором опыте используется коэффициент $\gamma=\gamma_2$
из~\eqref{eq:adaptation-gains}.
Результаты приведены на рисунках~\ref{fig:gamma2-error}
и~\ref{fig:gamma2-control}.

\imageplaceholder{lab12_gamma2_error.png}
  {Ошибка слежения $\varepsilon(t)$ при $\gamma=\gamma_2$}
  {fig:gamma2-error}

\imageplaceholder{lab12_gamma2_u.png}
  {Управляющее воздействие $u(t)$ при $\gamma=\gamma_2$}
  {fig:gamma2-control}

\subsection{Анализ результатов и сравнение опытов}

\reporttodo{сравнить два значения $\gamma$ по скорости затухания ошибки,
величине выбросов, остаточным колебаниям и амплитуде управления}

\noindent\textbf{Вывод по пункту:}
\reporttodo{сформулировать влияние коэффициента адаптации
и обосновать выбор предпочтительного значения для данного объекта}

\section{Выводы}

\reporttodo{кратко перечислить результаты синтеза обратной связи,
фильтра и адаптивного регулятора; оценить выполнение цели слежения
по графикам и сформулировать влияние $\gamma$}
